\section{Load Balancing}

\label{sec:loadbalancing}

\subsection{Adaptive Weight Computation}

The load balancer assigns weights to provider endpoints based on observed performance. Let $w_i(t)$ be the weight of provider $i$ at time $t$:

\begin{equation}
w_i(t) = \frac{H_i(t) \cdot C_i(t)}{\sum_j H_j(t) \cdot C_j(t)}
\end{equation}

where $H_i(t)$ is the health score and $C_i(t)$ is the remaining rate limit capacity estimated from provider-returned headers:

\begin{equation}
C_i(t) = \min\left(\frac{r_{\text{remaining},i}}{r_{\text{limit},i}}, \frac{t_{\text{remaining},i}}{t_{\text{limit},i}}\right)
\end{equation}

This formulation naturally routes traffic away from degraded or rate-limited providers without requiring explicit configuration.

\subsection{Fallback Chains}

Fallback chains define ordered sequences of provider endpoints to attempt for a given request. When the primary provider fails or is unavailable, the gateway transparently retries with the next provider in the chain:

\begin{algorithm}[t]
\caption{Fallback Chain Execution}
\label{alg:fallback}
\begin{algorithmic}[1]
\REQUIRE Request $q$, chain $[p_1, \ldots, p_n]$, retries $R$
\STATE $\text{errors} \gets []$
\FOR{$i = 1$ \TO $n$}
    \FOR{$r = 1$ \TO $R$}
        \STATE $\text{resp} \gets \text{dispatch}(q, p_i)$
        \IF{$\text{resp.success}$}
            \RETURN $\text{resp}$
        \ELSIF{$\text{resp.retryable}$}
            \STATE $\text{sleep}(\text{backoff}(r))$
            \STATE $\text{errors.append}(\text{resp.error})$
        \ELSE
            \STATE $\text{errors.append}(\text{resp.error})$
            \STATE \textbf{break} \COMMENT{Non-retryable, try next provider}
        \ENDIF
    \ENDFOR
\ENDFOR
\RETURN $\text{AggregateError}(\text{errors})$
\end{algorithmic}
\end{algorithm}

\subsubsection{Streaming Failover}

Streaming failover presents unique challenges. When a provider fails mid-stream, the gateway must decide whether to: (a) return the partial response, (b) retry from scratch with a different provider, or (c) attempt to continue from the partial output.

We implement strategy (b) with optimization: the gateway buffers streamed tokens and, upon failure, prepends them as an assistant prefix in the retry request. This approach works for providers supporting assistant prefilling (Anthropic, most open-weight model servers) and falls back to full retry otherwise.

\subsection{Priority Queues}

The gateway implements priority-based request scheduling with four levels:

\begin{enumerate}[leftmargin=*]
    \item \textbf{Critical}: Health checks, auth validation---bypass rate limits.
    \item \textbf{High}: Interactive user requests---minimize latency.
    \item \textbf{Normal}: API calls from applications---balance cost and latency.
    \item \textbf{Low}: Batch processing, embeddings---optimize for cost.
\end{enumerate}

Low-priority requests are eligible for delayed dispatch to exploit off-peak pricing windows where available.

